\chapter{Introdução}

A arquitetura de software é um campo de estudo que busca entender como sistemas de software são construídos e organizados para atender aos seus requisitos. Todavia, na prática, projetos costumam dedicar mais atenção às funcionalidades do que aos requisitos de qualidade. Isso traz à tona a necessidade de estudar e entender os padrões arquiteturais existentes e como eles impactam, a longo prazo, a confiabilidade, manutenibilidade, testabilidade, entre outras complicações para um sistema.

Este trabalho tem como objetivo estudar e entender um padrão arquitetural específico chamado de Modelo de Atores. O Modelo de Atores é originalmente uma teoria matemática largamente usada na arquitetura de sistemas distribuídos. Esse modelo enxerga sistemas computacionais em termos de unidades autônomas chamadas de atores que se comunicam entre si através de mensagens. Posteriormente, essa teoria serviu de paradigma base para a criação de linguagens de programação como Erlang e Elixir. 

Aqui é feita uma experimentação com o Modelo de Atores tentando aplicá-lo em um contexto distante do qual ele foi originalmente concebido. Isso é feito com a proposição de um arcabouço genérico para a implementação de aplicações centralizadas em Rust. Para validar o arcabouço, foi realizado um estudo de caso aplicando a proposta no projeto patch-hub, uma aplicação de terminal (ou seja, não distribuída)
\todo[inline]{Peguei a ideia, mas será que não vale deixar mais claro o que você ta querendo dizer?

Quem estiver lendo pode pensar ``Mas quando eu rodo `docker run' de uma imagem que não tá local, ele faz uma request e [...]'' (mesmo isto não sendo um critério para falar se uma aplicação é ou não distribuída, porque, senão, o patch-hub também seria o caso)

Uma ideia: e se você inverter e falar ``[...] uma aplicação não distribuída (uma interface de terminal para mantenedores do kernel Linux) [...]'' ou simplesmente falar direto ``[...] uma aplicação não distribuída [...]''.}
desenvolvida com a linguagem de programação Rust (que não possui o Modelo de Atores como paradigma central). A ideia é entender o esforço necessário para se aplicar um padrão arquitetural fora do seu contexto natural e quais consequências isso traz para o projeto.

Por que sequer seguir adiante com este estudo considerando a premissa de que o Modelo de Atores não é apropriado para o domínio desejado? Quando alguém procura um material resistente, certamente algodão e água não seriam as primeiras opções a serem pensadas. Todavia, ao se congelar um algodão encharcado, se obtém picrete, um material com resistência comparável ao concreto e que derrete muito mais lentamente que gelo convencional. Isso não é uma garantia de sucesso para combinações pouco usuais, mas significa que pode existir grande potencial a ser explorado.
\todo[inline]{Acho que nessa frase, você pode deixar mais concreto o que é o \textit{Isso} e trocar \textit{potencial a ser explorado} por \textit{potencial não explorado}. Algo como ``Esta ilustração não prova que combinações pouco usuais serão garantidamente bem-sucedidas, mas pontua como, às vezes, possuem um grande potencial não explorado.'' (desculpa se deturpei o sentido original que você queria; de toda forma, só uma sugestão!)}

De modo geral, é constatado que é preciso se depreender um esforço, tanto para se adaptar o arcabouço para um novo domínio, quanto para se reimaginar o problema sendo resolvido através de uma nova ótica e que, a partir disso, existe sim valor a ser extraído. Por meio do estudo de caso, pode-se concluir que não é impossível tal adaptação ser bem-sucedida e que talvez os padrões arquiteturais sejam mais flexíveis do que se imagina. O sucesso da adaptação está na preservação de uma parte considerável dos benefícios do padrão, mas isso depende da criatividade e do esforço do desenvolvedor.

\todo[inline]{Será que vale fazer um parágrafo curtinho aqui na linha de ``No contexto desta tese, o estudo de caso indicou que [...]''. Em outras palavras, um parágrafo ``respondendo'' os pontos que você colocou no parágrafo anterior no contexto do seu trabalho; ou seja, se concluiu que a adaptação no cenário montado foi bem sucedida, como/quanto do Actor Model foi preservado/conservado, se foi muito esforçoso, etc.}

O restante do texto está estruturado da seguinte forma: No capítulo 2, é feita uma explanação sobre o que é a arquitetura de software e o que são padrões arquiteturais; No capítulo 3, é apresentado o Modelo de Atores do ponto de vista teórico e prático com foco na linguagem de programação Elixir; No capítulo 4, é demonstrada a proposta de um arcabouço genérico para a implementação de aplicações centralizadas usando o Modelo de Atores em Rust com as devidas adaptações; No capítulo 5, é realizado um estudo de caso aplicando a proposta no projeto patch-hub; No capítulo 6, são apresentados os resultados da experimentação
\todo[inline]{Que tal ``[...] da experimentação usando métricas qualitativas e quantitativas do ponto de vista de engenharia de software''?}
; No capítulo 7, são apresentadas as conclusões finais e sugestões para futuros trabalhos.
